\section{Benchmark Models And Evaluation Framework}
\label{sec:benchmark_models_evaluation}

This section defines the benchmark protocol used to evaluate whether optical prescription data can be modeled by modern generative methods and whether optics-aware processing improves generated design quality. We evaluate unconditional generation because it isolates distribution learning: the model must generate plausible lens prescriptions and design specifications without being given a target condition at sampling time.

\subsection{Benchmark Models}\label{subsec:off_shelf_uncond_gen}

The benchmark suite uses mixed-type tabular generative models as the primary unconditional baselines, including STaSY, CoDi, and TabSyn \cite{Kim2023STaSy,Lee2023CoDi,Zhang2024TabSyn}. These methods are appropriate baselines because the model-facing lens representations are stored as mixed numerical and categorical tables. At the same time, they do not natively encode optical surface ordering, variable surface count, aperture-stop structure, or ray-trace validity constraints. This makes them useful for testing how much benefit comes from the proposed data representation and post-processing pipeline.

For each model family, experiments are organized across the dataset variants in Table~\ref{tab:dataset_variants_splits} and across the untuned and tuned model-facing representations. In the benchmark artifacts used for the reported results, the canonical training and validation partitions are combined for model fitting, while the test partition remains held out. Model-specific processors are allowed to perform required formatting, imputation, and normalization, but these steps are recorded separately from the common representation layer. This separation is important because the common data representation is intended to define the optical variables, while model-specific processors define the input format required by a particular generative model implementation.

\begin{table}[t]
\centering
\small
\caption{\textbf{Dataset variants used for benchmark training and evaluation.} Training counts report the rows provided to benchmark-model training, and test counts report the held-out rows used for evaluation. Counts include both patent-derived and synthetic rows.}
\label{tab:dataset_variants_splits}
\begin{tabular}{lrrr}
\toprule
\textbf{Variant} & \textbf{Training rows} & \textbf{Held-out test rows} & \textbf{Total rows} \\
\midrule
Patent-only & 1,128 & 125 & 1,253 \\
Patent+Synth-250k & 362,114 & 112,642 & 474,756 \\
Patent+Synth-500k & 612,114 & 112,642 & 724,756 \\
Patent+Synth-full & 1,013,691 & 112,642 & 1,126,333 \\
\bottomrule
\end{tabular}
\end{table}

% \begin{table}[t]
% \centering
% \footnotesize
% \caption{\textcolor{red}{[placeholder table]} \textbf{Benchmark training protocol.} This table will be filled with the final model list, dataset variants, representation settings, and sampling counts used in the benchmark runs.}
% \label{tab:benchmark_protocol}
% \begin{tabular}{p{0.18\linewidth}p{0.20\linewidth}p{0.25\linewidth}p{0.25\linewidth}}
% \toprule
% \textbf{Model} & \textbf{Representation(s)} & \textbf{Dataset variant(s)} & \textbf{Role} \\
% \midrule
% TabSyn & Tuned, untuned & Patent-only, Patent+Synth-250k, Patent+Synth-500k, Patent+Synth-full & Main latent-diffusion tabular baseline \\
% STaSy & Tuned, untuned & Planned benchmark variants & Score-based tabular baseline \\
% CoDi & Tuned, untuned & Planned benchmark variants & Mixed-type contrastive diffusion baseline \\
% SMOTE-style interpolation & Tuned or untuned reference & Selected variants & Non-neural diversity/memorization reference \\
% \bottomrule
% \end{tabular}
% \end{table}

\subsection{Evaluation Stages}

Evaluation is separated into stages so that local generated-prescription quality is not conflated with CodeV/LLESP simulation outcomes. Before CodeV results are available, generated samples can be evaluated for table validity, material validity, geometric plausibility, distribution matching, diversity, novelty, and LLESP-readiness after deterministic post-processing. After CodeV/LLESP result files return, the same samples can be evaluated for simulation error rate and optical performance metrics.

For unconditional generation, design conditions such as EFL, EPD, and maximum field angle are treated as generated design specifications. Their distributions are compared against reference data, but they are not treated as target-satisfaction metrics because no requested target is supplied to the model at sampling time.

\subsection{Metric Families}

Distribution and diversity metrics are computed on prescription and design-condition columns, excluding optical performance columns such as spot size and encircled-energy metrics. Optical performance is evaluated only after ray tracing or CodeV/LLESP simulation has produced valid result files; detailed post-processing attrition counts can be reported in the appendix if needed.

% Internal planning table retained for reference but omitted from the rendered paper.
\begin{comment}
\begin{table}[t]
\centering
\footnotesize
\caption{\textbf{Evaluation metric groups.} Metrics are separated by evaluation stage so that generated-prescription validity, distributional fidelity, simulation robustness, optical performance, and ray-tracing post-processing benefit are evaluated without conflating distinct failure modes.}
\label{tab:evaluation_metric_groups}
\begin{tabular}{>{\raggedright\arraybackslash}p{0.18\linewidth}>{\raggedright\arraybackslash}p{0.33\linewidth}>{\raggedright\arraybackslash}p{0.15\linewidth}>{\raggedright\arraybackslash}p{0.22\linewidth}}
\toprule
\textbf{Metric family} & \textbf{Primary quantities} & \textbf{Stage} & \textbf{Claim supported} \\
\midrule
Prescription distribution & Surface-count distribution distance, active numerical marginal KS, active \texttt{isair} marginal error, and pairwise correlation error within surface-count groups & Pre-CodeV & Fidelity to reference design distributions \\
Novelty / train proximity & Same-surface-count DCR excess relative to the expected train-closer probability from train/test support sizes & Pre-CodeV & Train-proximity check after correcting for split-size imbalance \\
Duplicate and leakage checks & Exact or rounded duplicate rates, generated rows matching training rows, and parent-family split checks & Pre-CodeV & Defensible separation between training data and generated samples \\
Prescription validity & Active-surface masks, material no-match rate, geometric intersections, semi-diameter sanity, and LLESP-readiness & Post-processing / pre-CodeV & Fraction of generated prescriptions usable before simulation \\
CodeV/LLESP outcomes & Submitted rows, returned rows, \texttt{!error} rate, and simulation success/failure categories & CodeV-dependent & Simulation robustness of generated prescriptions \\
Optical performance & SPO67, SPO100, EE50, EE90, and axial nominal performance across evaluated fields & CodeV-dependent & Optical quality after simulation \\
Ray-tracing correction benefit & Paired feasible-fraction and optical-performance changes with and without physics-based post-processing & CodeV-dependent paired ablation & Effect of physics-based post-processing \\
\bottomrule
\end{tabular}
\end{table}
\end{comment}

For novelty and train-proximity diagnostics, nearest-neighbor comparisons are restricted to the same surface-count group. This avoids treating padded inactive surface slots as ordinary features and is equivalent to assigning cross-surface-count pairs a larger topology-aware distance than any same-surface-count pair. The reported DCR excess subtracts the expected train-closer probability implied by the same-surface-count train/test support sizes, so it is not interpreted against a fixed 50\% baseline unless the train and test supports are equal.
